\documentclass[11pt]{article}

\usepackage[utf8]{inputenc}
\usepackage[T1]{fontenc}
\usepackage{lmodern}
\usepackage{geometry}
\usepackage{amsmath,amssymb,amsfonts,amsthm,mathtools}
\usepackage{bm}
\usepackage{enumitem}
\usepackage{microtype}
\usepackage{hyperref}
\usepackage{titlesec}
\usepackage{booktabs}
\usepackage{array}
\usepackage{setspace}
\usepackage{mathrsfs}
\usepackage{physics}

\geometry{margin=1in}
\hypersetup{colorlinks=true, linkcolor=black, urlcolor=blue, citecolor=black}
\setlist[enumerate]{topsep=0.4em,itemsep=0.35em}
\setlist[itemize]{topsep=0.3em,itemsep=0.25em}
\titleformat{\section}{\large\bfseries}{\thesection.}{0.5em}{}
\titleformat{\subsection}{\normalsize\bfseries}{\thesubsection.}{0.5em}{}
\setstretch{1.06}

\newcommand{\Aether}{\AE{}ther}
\newcommand{\Aflow}{\AE{}ther-flow}
\newcommand{\TheoryName}{\Aflow{} Interpretation of Relativity}
\newcommand{\TheoryProgram}{\Aether{} / \Aflow{} framework}
\newcommand{\CompletedMetric}{g^{\sharp}}
\newcommand{\CompletionCore}{\mathfrak C^{\mathrm{zr,aff}}}
\newcommand{\TraceCore}{\mathfrak C^{\mathrm{tr}}}
\newcommand{\AffineNucleus}{\mathfrak N^{\mathrm{aff}}}
\newcommand{\AmbientPackage}{\mathfrak A^{\mathrm{bal}}}
\newcommand{\CombinedLaw}{\mathfrak R^{\mathrm{comb}}}
\newcommand{\SelectorLaw}{\mathfrak S^{\mathrm{sel},0}}

\newtheorem{definition}{Definition}
\newtheorem{proposition}{Proposition}
\newtheorem{remark}{Remark}
\newtheorem{corollary}{Corollary}
\newtheorem{theorem}{Theorem}

\title{\textbf{The \TheoryName{}}\\[0.4em]
\large Primitive-Reservoir Observer-Reduction\\
\large Zero-Hidden Selector Generator Law\\
\large Next Benchmark Criterion}
\author{}
\date{}

\begin{document}

\maketitle

\begin{abstract}
The current combined-response frontier note had already fixed the next honest burden: go one level below the completion-balance ambient dynamics theorem and try to derive or uniquely force the combined response substrate law \(\CombinedLaw\) itself from still more primitive explicit substrate structure, with fallback to a smaller theorem inside that law if the full forcing could not yet be done cleanly. The new selector-law forcing theorem has now spent that preferred option: the combined response substrate law is itself a canonical and necessary reduction of a deeper zero-hidden selector generator law \(\SelectorLaw\).

This note records the routing consequence of that new theorem. The next honest benchmark-facing manuscript must therefore go one level below the current selector frontier: first derive or uniquely force the zero-hidden selector generator law \(\SelectorLaw\) from still more primitive explicit substrate structure in a way that actually changes benchmark-facing status, or, if that cannot be done cleanly, prove a genuinely smaller theorem-level collapse, sharper necessity result, or sharper uniqueness result strictly inside that selector law. Another ambient-package restatement, combined-response restatement, trace-core restatement, completion-core restatement, exact-shell affine nucleus restatement, or another selector-law relay that preserves the same selector law without shrinking or forcing it further does not count next.

The benchmark boundary remains fixed throughout. The one operative metric is still exactly \(\CompletedMetric\), the ordinary matter route is unchanged, there is no second operative metric, the low-energy matter law is not enlarged, and stronger first-principles promotion language remains paused. The positive result of this note is therefore routing discipline at the new selector-law frontier.
\end{abstract}

\section{Introduction}

The active derivational program of \TheoryName{} is no longer at a stage where the combined response substrate law \(\CombinedLaw\) should be treated as the default unexplained stopping point. The current ambient-dynamics frontier had already shown that the ambient package \(\AmbientPackage\) is forced by that deeper combined response law. The new selector-law forcing theorem then spent the preferred next move available there: \(\CombinedLaw\) itself is now forced by the deeper zero-hidden selector generator law \(\SelectorLaw\).

The present note records the routing consequence of that theorem. It does not reopen the selector-law mathematics, and it does not claim a completed first-principles substrate derivation of gravity. It records only which kind of next manuscript would now count as an honest benchmark-facing gain below the current selector frontier.

\paragraph{Framework and claim boundary.}
\TheoryProgram{} denotes the broader \Aether{} / \Aflow{} framework built on the claims that the \Aether{} is the underlying four-dimensional substrate of reality, the \Aflow{} is its intrinsic ordered motion, observed three-dimensional space is the local experiential slice of that deeper substrate, S-time is the experienced order of change arising from matter, light, and the \Aflow{}, and observed expansion is the three-dimensional appearance of deeper four-dimensional ordered motion. Within that framework, \TheoryName{} denotes the exact-closure theory in which the effective gravitational dynamics are adopted to be exactly Einsteinian with universal matter coupling. In the active sequence, \emph{adoption} means use of that established relativistic dynamics without claiming substrate derivation, while \emph{derivation} is reserved for a first-principles recovery from explicit substrate variables.

The present note stays strictly inside that boundary. It records only which kind of next manuscript would now count as an honest benchmark-facing gain below the selector-law frontier.

\section{Fixed Inputs}

\begin{definition}[Selector-frontier fixed inputs]
The criterion recorded here uses the following fixed inputs.
\begin{enumerate}
    \item The benchmark exact-closure package, which fixes the public one-metric GR target of the repository.
    \item The benchmark gatekeeping layer, including the post-bridge status correction and the upstream-compression safeguard that blocks same-output deeper-origin relays from counting as new front-line progress once the benchmark-relevant object is unchanged.
    \item The current ambient-dynamics forcing theorem, which proves that the completion-balance ambient package \(\AmbientPackage\) is a canonical and necessary reduction of the deeper combined response substrate law \(\CombinedLaw\).
    \item The new selector-law forcing theorem, which proves that the combined response substrate law \(\CombinedLaw\) is itself a canonical and necessary reduction of the deeper zero-hidden selector generator law \(\SelectorLaw\).
\end{enumerate}
\end{definition}

\begin{remark}
The present note does not replace those manuscripts. It records the routing consequence of reading them together under the active safeguard against recursive depth drift.
\end{remark}

\section{Why the Selector-Law Theorem Is the Frontier}

\begin{definition}[Benchmark-neutral selector relay]
A \emph{benchmark-neutral selector relay} is a manuscript that preserves the same benchmark one-metric output, the same ordinary matter route, the same exact-shell affine nucleus consequence, the same trace-core consequence, the same ambient-package consequence, the same combined-response consequence, and the same claim boundary while merely restating the ambient package \(\AmbientPackage\), the combined response law \(\CombinedLaw\), the trace core \(\TraceCore\), the full completion-core presentation \(\CompletionCore\), the exact-shell affine nucleus \(\AffineNucleus\), or the same zero-hidden selector generator law \(\SelectorLaw\) without shrinking or forcing that selector law further.
\end{definition}

\begin{proposition}[The selector-law forcing theorem is the live frontier]
At the current stage of the primitive-reservoir observer-reduction line, the theorem deriving the combined response substrate law \(\CombinedLaw\) from the zero-hidden selector generator law \(\SelectorLaw\) is the live benchmark-facing frontier.
\end{proposition}

\begin{proof}
That theorem changes benchmark-facing status at current scope because it removes the combined response law \(\CombinedLaw\) itself as the live internal stopping point. The current bridge datum is no longer merely a deeper combined-response package whose ambient consequence was already known; it is now a canonical output of the deeper selector law \(\SelectorLaw\). By contrast, a manuscript that preserves the same benchmark package, the same combined-response consequence, and the same selector-law stopping point while merely rephrasing the ambient package, combined response law, trace core, completion core, or exact-shell affine nucleus does not alter benchmark-facing status unless it adds a comparably sharp gain. Therefore the live frontier is the selector-law forcing theorem rather than the earlier combined-response ambient theorem or a later same-output selector relay.
\end{proof}

\begin{corollary}[Status of the combined-response law as a next-step target]
Under the current routing safeguard, the combined response substrate law \(\CombinedLaw\) is no longer itself the default next benchmark-facing target.
\end{corollary}

\begin{proof}
The selector-law forcing theorem already proves that \(\CombinedLaw\) is forced by \(\SelectorLaw\). Once that is on record, another manuscript that spends the move on the combined-response law again does not remove any further benchmark-relevant freedom. The open burden has therefore moved below the combined-response law itself.
\end{proof}

\section{The Next Benchmark Criterion}

\begin{definition}[Admissible next benchmark-facing selector manuscript]
At the current selector frontier, a manuscript counts as an admissible next benchmark-facing move only if it respects the following burden order:
\begin{enumerate}
    \item its primary attempt is to derive or uniquely force the zero-hidden selector generator law \(\SelectorLaw\) from still more primitive explicit substrate structure in a way that does more than merely relocate the unexplained layer deeper;
    \item if that cannot be done cleanly at current scope, its admissible fallback gain is to prove a genuinely smaller theorem-level collapse, sharper necessity result, or sharper uniqueness result strictly inside \(\SelectorLaw\);
    \item in either case it preserves the same benchmark one-metric package, the same ordinary matter route, and the same claim boundary, and it introduces neither a second operative metric nor an enlarged low-energy matter law and licenses no stronger first-principles promotion language.
\end{enumerate}
\end{definition}

\begin{theorem}[Next honest benchmark-facing task at selector scope]
The next honest benchmark-facing manuscript at the current selector frontier should therefore go one level below that frontier: first try to derive or uniquely force the zero-hidden selector generator law \(\SelectorLaw\) from still more primitive explicit substrate structure in a way that actually changes benchmark-facing status. If that cannot be done cleanly, the admissible fallback is to prove a genuinely smaller collapse, sharper necessity result, or sharper uniqueness result strictly inside that selector law. A manuscript that only restates the completion-balance ambient package \(\AmbientPackage\), only restates the combined response law \(\CombinedLaw\), only restates the trace core \(\TraceCore\), only restates the full completion-core presentation \(\CompletionCore\), only restates the exact-shell affine nucleus \(\AffineNucleus\), or only repeats the same selector-law consequence while preserving the same selector law without shrinking or forcing it further does not satisfy the criterion.
\end{theorem}

\begin{proof}
By Proposition 1, the live frontier already lies at the theorem that forces the current combined response law from the deeper selector law. Therefore a second manuscript below it must improve on that status rather than merely accompany it. A deeper-origin manuscript can do so only if it changes benchmark-facing status by more than depth alone, and the first clean way to do that is to force \(\SelectorLaw\) itself from still more primitive explicit substrate structure. If that first option is not yet available cleanly, the routing safeguard allows the fallback of a smaller internal collapse, necessity result, or uniqueness result inside \(\SelectorLaw\). If a manuscript simply reproduces the same selector-law combined-response consequence through another relay while preserving the same selector law without shrinking or forcing it further, the routing rule classifies it as benchmark neutral.
\end{proof}

\section{What the Next Move Should Not Be}

\begin{proposition}[Screened next moves]
At the current selector frontier, none of the following is the default next benchmark-facing move:
\begin{enumerate}
    \item another restatement of the completion-balance ambient dynamics package \(\AmbientPackage\);
    \item another restatement of the combined response substrate law \(\CombinedLaw\);
    \item another restatement of the completion-state shell-trace core \(\TraceCore\);
    \item another restatement of the full completion-core graph presentation \(\CompletionCore\);
    \item another restatement of the exact-shell affine nucleus \(\AffineNucleus\);
    \item another selector-law relay that preserves the same zero-hidden selector generator law without shrinking or forcing it further;
    \item a manuscript that introduces a second operative metric, enlarges the ordinary matter law, or uses stronger first-principles promotion language.
\end{enumerate}
\end{proposition}

\begin{proof}
Items (1)--(6) fail the preceding criterion because they do not directly address the current selector-law burden. They revisit either already-forced smaller bridge objects, already-forced larger bridge objects, or the same selector-law stopping point rather than removing the remaining selector-law freedom or proving a smaller internal datum inside \(\SelectorLaw\). Item (7) violates the fixed claim boundary of the benchmark package and therefore cannot count as a disciplined next step on the active line. The benchmark package remains the exact one-metric GR package already fixed by adoption, so any admissible next manuscript must preserve that boundary while keeping stronger first-principles promotion language paused.
\end{proof}

\begin{remark}
This screening statement is not a denial that the ambient-package theorem, the combined-response theorem, the trace-core theorem, the completion-core theorem, or older relays matter historically. It is a routing statement: they are not the honest way to spend the next benchmark-facing move from the current selector frontier.
\end{remark}

\section{Conclusion}

The positive result of this note is routing discipline at the current selector frontier. The selector-law forcing theorem remains the live benchmark-facing gain because it already removes the combined response law as the internal stopping point on the active line. The combined response law therefore no longer sets the honest next burden, and neither do the already-forced ambient package, trace core, completion-core presentation, or exact-shell affine nucleus.

The scope remains narrow and honest. The benchmark package is unchanged: the one operative metric is still exactly \(\CompletedMetric\), the ordinary matter route is unchanged, there is no second operative metric, and the low-energy matter law is not enlarged. Nothing here upgrades adoption to derivation or licenses stronger first-principles promotion language yet.

The next open burden is therefore clear. The next benchmark-facing manuscript should go one level below the current selector frontier: first try to derive or uniquely force the zero-hidden selector generator law \(\SelectorLaw\) from still more primitive explicit substrate structure in a way that actually changes benchmark-facing status while preserving the same benchmark one-metric package and claim boundary. If that cannot be done cleanly, the admissible fallback is a genuinely smaller collapse, sharper necessity result, or sharper uniqueness result strictly inside that selector law. Another ambient-package restatement, combined-response restatement, trace-core restatement, completion-core restatement, exact-shell affine nucleus restatement, or another selector-law relay that preserves the same selector law without shrinking or forcing it further is not the clean next manuscript target at current scope.

\end{document}
